% Generated draft: requires mathematical and source review.
\documentclass[11pt]{article}
\usepackage[margin=25mm]{geometry}
\usepackage{fontspec}
\setmainfont{Latin Modern Roman}
\usepackage{amsmath,amssymb,mathtools,unicode-math}
\setmathfont{Latin Modern Math}
\usepackage{enumitem,needspace,xurl,hyperref}
\setlength{\parindent}{0pt}
\setlength{\parskip}{5pt}
\setlength{\emergencystretch}{3em}
\newcommand{\N}{\mathbb{N}}
\newcommand{\Z}{\mathbb{Z}}
\newcommand{\Q}{\mathbb{Q}}
\newcommand{\R}{\mathbb{R}}
\newcommand{\C}{\mathbb{C}}
\newcommand{\F}{\mathbb{F}}
\newcommand{\A}{\mathbb{A}}
\newcommand{\PP}{\mathbb{P}}
\newcommand{\E}{\mathbb{E}}
\newcommand{\eps}{\varepsilon}
\newcommand{\dd}{\,\mathrm{d}}
\begin{document}
\section*{UnsolvedMath ID: 20003144}
\textbf{Source problem code:} catalogue entry
\subsection{Definitions and mathematical statement}
The source asks whether an $n$-torus can collapse to an interval and already records the answer ``No,'' essentially settled at the workshop. The suppressed setting, made explicit by the supplied verification, is a sequence of closed smooth Riemannian tori $(T^n,g_i)$ with one uniform lower sectional-curvature bound $\sec_{g_i}\geq\kappa$, for fixed $n$, converging in Gromov--Hausdorff distance to a nondegenerate interval $[0,L]$, $L>0$. The reported theorem rules out such a sequence. No two-sided curvature bound is required in this statement. A point is not the nondegenerate interval in question.

The aid credits Zamora with the subsequent general proof and reports later full torus stability for compact limits by Bru\`e--Naber--Semola, without giving precise hypotheses or a full statement of that stronger theorem. The label solved is a supplied literature assessment; the detailed mathematical result in this particular attempt is a restricted cohomogeneity-one model obstruction.

For that obstruction let a connected closed $(k+1)$-manifold $M$, $k\geq1$, have an interval-type cohomogeneity-one $T^k$ group diagram with finite principal isotropy $H$, and endpoint isotropy groups $K_-,K_+$ satisfying $K_\pm/H\cong S^1$. The two disk-bundle pieces glue along the principal orbit $T^k/H$. Using this common orbit to identify rational first homology with $\Q^k$, let $L_-,L_+\subset\Q^k$ be the nonzero one-dimensional rational slope subspaces killed by the endpoint disk attachments. They must be transported through the chosen equivariant gluing before comparison. The aid's formula is
\[
 H_1(M;\Q)\cong\Q^k/(L_-+L_+),
\]
\[
 b_1(M;\Q)=k-2+\dim_{\Q}(L_-\cap L_+)\leq k-1.
\]
The intersection dimension is one for equal slope lines and zero otherwise. In particular this total space cannot be $T^{k+1}$, whose rational first Betti number is $k+1$. Finite isotropy may matter integrally; this assertion uses rational homology.

This excludes the specified global torus-action diagrams. It does not establish that every arbitrary lower-curvature collapsing sequence has such a global cohomogeneity-one model, and therefore is not itself the full proof of the general no-interval theorem.
\subsection{Short English statement}
Closed Riemannian tori with a uniform lower sectional-curvature bound cannot converge to a nondegenerate interval, according to the supplied historical answer and later literature assessment. The attempt proves a more restricted model obstruction: an interval-type torus action with circle endpoint isotropy kills too much rational first homology for its total space to be a torus.
\subsection{Sources}
\begin{itemize}
\item Local UnsolvedMath registry, permanent ID $20003144$, source code catalogue entry; canonical location \texttt{aim-topology-notes.json}, \texttt{notes[231]}.
\item American Institute of Mathematics, \emph{Manifolds with nonnegative sectional curvature}, Problem 19, \url{https://aimath.org/WWN/nnsectcurvature/nnsectcurvature.pdf}. The original sentence explicitly says the answer is ``No,'' essentially settled at the workshop. The supplied verification specifies the uniform lower-curvature GH-collapse setting.
\item Zamora, and later Bru\`e--Naber--Semola, as credited in the attached research summary; no titles, dates or full bibliographic references are supplied.
\item Research attempt 2 is marked partial result, proved, and low novelty confidence, while the registry labels the historical question solved in the literature. The explicit caveat identifies the aid as machine generated, not peer reviewed, and requiring independent verification before citation.
\end{itemize}
\end{document}
